\newglossaryentry{drawing-test}{
    name={Рисуночный тест},
    description={проективная психодиагностическая методика, основанная на интерпретации создаваемого испытуемым рисунка}
}

\newglossaryentry{cnn}{
    name={Свёрточная нейронная сеть},
    description={специальный класс искусственных нейронных сетей, использующий операцию свёртки для обработки данных с пространственной структурой, в первую очередь изображений}
}

\newglossaryentry{multi-task}{
    name={Multi-task learning},
    description={парадигма обучения, в которой одна модель одновременно решает несколько связанных задач, используя общее представление признаков}
}

\newglossaryentry{f1}{
    name={F1-мера},
    description={гармоническое среднее точности (precision) и полноты (recall) бинарного или мультиклассового классификатора}
}

\newglossaryentry{bbox}{
    name={Ограничивающий прямоугольник},
    description={прямоугольная область изображения, в которую заключён обнаруженный объект; описывается координатами левого верхнего и правого нижнего углов}
}

\newglossaryentry{adaptive-threshold}{
    name={Адаптивная бинаризация},
    description={метод бинаризации изображения, в котором порог вычисляется индивидуально для каждой точки на основе локальной окрестности}
}

\newglossaryentry{morphology}{
    name={Морфологические операции},
    description={базовые операции обработки бинарных изображений: эрозия, дилатация, замыкание (closing), размыкание (opening), основанные на структурирующем элементе}
}

\newglossaryentry{cc-labeling}{
    name={Выделение связных компонент},
    description={метод обработки бинарного изображения, при котором смежные единичные пиксели группируются в связные области; для каждой области вычисляются её координаты, размеры и площадь}
}
